\documentclass[10pt,a4paper]{amsart}
\usepackage{amsmath,amssymb,amsthm,graphicx}
% Retained minimal preamble for the Stacks Project chapter flat.tex (More on Flatness).
% Verbatim excerpt of the chapter preamble: Xy-pic setup (native preamble lines 14-18) and the native
% theorem declarations with the mathematical macros (native preamble lines 154-217). Copied mechanically;
% nothing added, removed or redefined. The wrapper class and the amsmath/amssymb/amsthm/graphicx base
% packages are supplied by the assembler, so the class/externaldocument part of the native preamble is omitted.
% For commutative diagrams we use Xy-pic
\usepackage[all]{xy}

% We use 2cell for 2-commutative diagrams.
\xyoption{2cell}
% Theorem environments.
%
\theoremstyle{plain}
\newtheorem{theorem}[subsection]{Theorem}
\newtheorem{proposition}[subsection]{Proposition}
\newtheorem{lemma}[subsection]{Lemma}

\theoremstyle{definition}
\newtheorem{definition}[subsection]{Definition}
\newtheorem{example}[subsection]{Example}
\newtheorem{exercise}[subsection]{Exercise}
\newtheorem{situation}[subsection]{Situation}

\theoremstyle{remark}
\newtheorem{remark}[subsection]{Remark}
\newtheorem{remarks}[subsection]{Remarks}

\numberwithin{equation}{subsection}

% Macros
%
\def\lim{\mathop{\mathrm{lim}}\nolimits}
\def\colim{\mathop{\mathrm{colim}}\nolimits}
\def\Spec{\mathop{\mathrm{Spec}}}
\def\Hom{\mathop{\mathrm{Hom}}\nolimits}
\def\Ext{\mathop{\mathrm{Ext}}\nolimits}
\def\SheafHom{\mathop{\mathcal{H}\!\mathit{om}}\nolimits}
\def\SheafExt{\mathop{\mathcal{E}\!\mathit{xt}}\nolimits}
\def\Sch{\mathit{Sch}}
\def\Mor{\mathop{\mathrm{Mor}}\nolimits}
\def\Ob{\mathop{\mathrm{Ob}}\nolimits}
\def\Sh{\mathop{\mathit{Sh}}\nolimits}
\def\NL{\mathop{N\!L}\nolimits}
\def\CH{\mathop{\mathrm{CH}}\nolimits}
\def\proetale{{pro\text{-}\acute{e}tale}}
\def\etale{{\acute{e}tale}}
\def\QCoh{\mathit{QCoh}}
\def\Ker{\mathop{\mathrm{Ker}}}
\def\Im{\mathop{\mathrm{Im}}}
\def\Coker{\mathop{\mathrm{Coker}}}
\def\Coim{\mathop{\mathrm{Coim}}}

% Boxtimes
%
\DeclareMathSymbol{\boxtimes}{\mathbin}{AMSa}{"02}

%
% Macros for moduli stacks/spaces
%
\def\QCohstack{\mathcal{QC}\!\mathit{oh}}
\def\Cohstack{\mathcal{C}\!\mathit{oh}}
\def\Spacesstack{\mathcal{S}\!\mathit{paces}}
\def\Quotfunctor{\mathrm{Quot}}
\def\Hilbfunctor{\mathrm{Hilb}}
\def\Curvesstack{\mathcal{C}\!\mathit{urves}}
\def\Polarizedstack{\mathcal{P}\!\mathit{olarized}}
\def\Complexesstack{\mathcal{C}\!\mathit{omplexes}}
% \Pic is the operator that assigns to X its picard group, usage \Pic(X)
% \Picardstack_{X/B} denotes the Picard stack of X over B
% \Picardfunctor_{X/B} denotes the Picard functor of X over B
\def\Pic{\mathop{\mathrm{Pic}}\nolimits}
\def\Picardstack{\mathcal{P}\!\mathit{ic}}
\def\Picardfunctor{\mathrm{Pic}}
\def\Deformationcategory{\mathcal{D}\!\mathit{ef}}
\usepackage{cleveref}
\crefname{originallemma}{original lemma}{original lemmas}
\Crefname{originallemma}{Original lemma}{Original lemmas}
\makeatletter
\newlabel{coherent-lemma-affine-base-change}{{02KG}{source}{}{sourcepointer-1}{}}
\newlabel{coherent-lemma-affine-base-change@cref}{{[originallemma][0][]02KG}{[1][0][]source}{}{}{}}
\makeatother
\crefname{originallemma}{original lemma}{original lemmas}
\Crefname{originallemma}{Original lemma}{Original lemmas}
\makeatletter
\newlabel{descent-lemma-dimension-at-point-local}{{04N4}{source}{}{sourcepointer-2}{}}
\newlabel{descent-lemma-dimension-at-point-local@cref}{{[originallemma][0][]04N4}{[1][0][]source}{}{}{}}
\makeatother
\crefname{originallemma}{original lemma}{original lemmas}
\Crefname{originallemma}{Original lemma}{Original lemmas}
\makeatletter
\newlabel{lemma-lift-etale}{{057R}{source}{}{sourcepointer-3}{}}
\newlabel{lemma-lift-etale@cref}{{[originallemma][0][]057R}{[1][0][]source}{}{}{}}
\makeatother
\crefname{originallemma}{original lemma}{original lemmas}
\Crefname{originallemma}{Original lemma}{Original lemmas}
\makeatletter
\newlabel{more-morphisms-lemma-normal-morphism-irreducible}{{057J}{source}{}{sourcepointer-4}{}}
\newlabel{more-morphisms-lemma-normal-morphism-irreducible@cref}{{[originallemma][0][]057J}{[1][0][]source}{}{}{}}
\makeatother
\crefname{originallemma}{original lemma}{original lemmas}
\Crefname{originallemma}{Original lemma}{Original lemmas}
\makeatletter
\newlabel{varieties-lemma-geometrically-connected-if-connected-and-point}{{04KV}{source}{}{sourcepointer-5}{}}
\newlabel{varieties-lemma-geometrically-connected-if-connected-and-point@cref}{{[originallemma][0][]04KV}{[1][0][]source}{}{}{}}
\makeatother
\title{Refining a one step devissage along an elementary etale neighbourhood of the base}
\author{The Stacks Project Authors}
\date{stacks-project:flat.tex:local-snapshot-20261011; GNU FDL 1.2 or later}
\begin{document}
\maketitle
\noindent\textit{Editorial note.} All delivered mathematical text is the author's own, copied from the pinned Stacks Project chapter file flat.tex; no statement or proof step was rewritten, reordered or removed. The sentence 'Of course in general this is not going to be true' and the surrounding argument at proof lines 1116-1134 are the author's own account of the only obstruction.\par\medskip
\noindent\textit{Editorial note.} The two leading context spans are the author's own definitions of a one step devissage over s and at x, printed first because the counted statement opens with a verbatim reference to the second of them; they define the labels definition-one-step-devissage and definition-one-step-devissage-at-x, which therefore do not appear among the external references.\par\medskip
\noindent\textit{Editorial note.} The preamble is the chapter's own retained minimal preamble (sources/stacks-flat/retained-preamble.tex), a verbatim excerpt of the native Xy-pic setup and of the native theorem declarations and mathematical macros. The native class line, the xr-hyper externaldocument block and the reference/slogan/history display environments are omitted because the assembler supplies a plain amsart wrapper and the excerpt uses none of those environments; no macro was invented or redefined.\par\medskip
\noindent\textit{Editorial note.} Locators of the cross-chapter references point either at the \textbackslash{}label line of the referenced result inside the pinned chapter file (same-chapter reference 057R) or at a verbatim excerpt of the project's own tag list, sources/stacks-flat/referenced-tags.txt, each recorded with its stable Stacks tag.\par\medskip
\noindent\textit{Allowed standard prerequisites (editorial).}\begin{itemize}
\item Elementary etale neighbourhoods of pointed schemes. The statement is about an elementary etale neighbourhood (Y', y') -> (Y, y) -> (S, s) and produces an elementary etale neighbourhood (X'', x'') -> (X, x); the construction of the new diagram is the retained obligation.
\item Invariance of the dimension of a support under etale localization. Used to keep dim\_\{x''\} Supp(F''\_\{s''\}) equal to n; the surrounding verifications of the definition are retained.
\item Existence of an elementary etale neighbourhood making all fibres geometrically irreducible. Supplies the final (S'', s'') -> (S', s') and affine Y'' once the proof itself has produced the smooth geometrically connected component T; the reduction to that component and the residue-field argument are retained.
\end{itemize}\medskip
\ifdefined\hypertarget\hypertarget{sourcepointer-1}{}\fi
\noindent\textit{Source reference.} Original lemma 02KG: sources/stacks-flat/referenced-tags.txt, source lines 5--5. Cross-chapter pointer used at proof lines 1084-1087 to identify the pushforward of G'' with F'' over affine schemes; it records an identification the counted proof itself sets up by construction.\par\medskip
\ifdefined\hypertarget\hypertarget{sourcepointer-2}{}\fi
\noindent\textit{Source reference.} Original lemma 04N4: sources/stacks-flat/referenced-tags.txt, source lines 6--6. Cross-chapter standard prerequisite (dimension at a point is etale local) invoked at proof lines 1099-1107 to keep the fibre support dimension equal to n; the residue-field and etale comparison argument around it is retained.\par\medskip
\ifdefined\hypertarget\hypertarget{sourcepointer-3}{}\fi
\noindent\textit{Source reference.} Original lemma 057R: sources/stacks-flat/pinned-original-source.tex, source lines 84--84. Same-chapter pointer to the etale lifting lemma invoked at proof lines 1072-1074 to build X'' with Z'' = X'' x\_\{X\_\{S'\}\} Z\_\{S'\}; the choice of Y'', S'' and the verification of all the devissage conditions remain in the counted proof.\par\medskip
\ifdefined\hypertarget\hypertarget{sourcepointer-4}{}\fi
\noindent\textit{Source reference.} Original lemma 057J: sources/stacks-flat/referenced-tags.txt, source lines 7--7. Cross-chapter input invoked at proof lines 1127-1132 for the final choice of (S'', s'') and Y''; the proof's own reduction to the smooth component T and its residue-field argument are retained and are what make the pointer applicable.\par\medskip
\ifdefined\hypertarget\hypertarget{sourcepointer-5}{}\fi
\noindent\textit{Source reference.} Original lemma 04KV: sources/stacks-flat/referenced-tags.txt, source lines 10--10. Cross-chapter criterion invoked at proof lines 1122-1125 to convert kappa(s) algebraically closed in kappa(y) into geometric connectedness of the component T; a navigation pointer for a standard criterion, not a replacement of the counted construction.\par\medskip
\begin{definition}
\label{definition-one-step-devissage}
Let $S$ be a scheme.
Let $X$ be locally of finite type over $S$.
Let $\mathcal{F}$ be a quasi-coherent $\mathcal{O}_X$-module of finite type.
Let $s \in S$ be a point.
A {\it one step d\'evissage of $\mathcal{F}/X/S$ over $s$}
is given by morphisms of schemes over $S$
$$
\xymatrix{
X & Z \ar[l]_i \ar[r]^\pi & Y
}
$$
and a quasi-coherent $\mathcal{O}_Z$-module $\mathcal{G}$ of finite type
such that
\begin{enumerate}
\item $X$, $S$, $Z$ and $Y$ are affine,
\item $i$ is a closed immersion of finite presentation,
\item $\mathcal{F} \cong i_*\mathcal{G}$,
\item $\pi$ is finite, and
\item the structure morphism $Y \to S$ is smooth with
geometrically irreducible fibres of
dimension $\dim(\text{Supp}(\mathcal{F}_s))$.
\end{enumerate}
In this case we say $(Z, Y, i, \pi, \mathcal{G})$ is a one step
d\'evissage of $\mathcal{F}/X/S$ over $s$.
\end{definition}

\begin{definition}
\label{definition-one-step-devissage-at-x}
Let $S$ be a scheme.
Let $X$ be locally of finite type over $S$.
Let $\mathcal{F}$ be a quasi-coherent $\mathcal{O}_X$-module of finite type.
Let $x \in X$ be a point with image $s$ in $S$.
A {\it one step d\'evissage of $\mathcal{F}/X/S$ at $x$}
is a system $(Z, Y, i, \pi, \mathcal{G}, z, y)$, where
$(Z, Y, i, \pi, \mathcal{G})$ is a one step d\'evissage of
$\mathcal{F}/X/S$ over $s$ and
\begin{enumerate}
\item $\dim_x(\text{Supp}(\mathcal{F}_s)) = \dim(\text{Supp}(\mathcal{F}_s))$,
\item $z \in Z$ is a point with $i(z) = x$ and $\pi(z) = y$,
\item we have $\pi^{-1}(\{y\}) = \{z\}$,
\item the extension $\kappa(y)/\kappa(s)$ is purely
transcendental.
\end{enumerate}
\end{definition}

\begin{lemma}
\label{lemma-elementary-etale-neighbourhood}
Let $S$, $X$, $\mathcal{F}$, $x$, $s$ be as in
Definition \ref{definition-one-step-devissage-at-x}.
Let $(Z, Y, i, \pi, \mathcal{G}, z, y)$ be a one step d\'evissage
of $\mathcal{F}/X/S$ at $x$. Let
$$
\xymatrix{
(Y, y) \ar[d] & (Y', y') \ar[l] \ar[d] \\
(S, s) & (S', s') \ar[l]
}
$$
be a commutative diagram of pointed schemes such that the horizontal
arrows are elementary \'etale neighbourhoods. Then there exists
a commutative diagram
$$
\xymatrix{
& & (X'', x'') \ar[lld] \ar[d] & (Z'', z'') \ar[l] \ar[lld] \ar[d] \\
(X, x) \ar[d] & (Z, z) \ar[l] \ar[d] &
(S'', s'') \ar[lld] & (Y'', y'') \ar[lld] \ar[l] \\
(S, s) & (Y, y) \ar[l]
}
$$
of pointed schemes with the following properties:
\begin{enumerate}
\item $(S'', s'') \to (S', s')$ is an elementary \'etale neighbourhood and
the morphism $S'' \to S$ is the composition $S'' \to S' \to S$,
\item $Y''$ is an open subscheme of $Y' \times_{S'} S''$,
\item $Z'' = Z \times_Y Y''$,
\item $(X'', x'') \to (X, x)$ is an elementary \'etale neighbourhood, and
\item $(Z'', Y'', i'', \pi'', \mathcal{G}'', z'', y'')$ is a one step
d\'evissage at $x''$ of the sheaf $\mathcal{F}''$.
\end{enumerate}
Here $\mathcal{F}''$ (resp.\ $\mathcal{G}''$) is the pullback of
$\mathcal{F}$ (resp.\ $\mathcal{G}$) via the morphism $X'' \to X$
(resp.\ $Z'' \to Z$) and $i'' : Z'' \to X''$ and $\pi'' : Z'' \to Y''$
are as in the diagram.
\end{lemma}

\begin{proof}
Let $(S'', s'') \to (S', s')$ be any elementary \'etale neighbourhood
with $S''$ affine. Let $Y'' \subset Y' \times_{S'} S''$ be any affine
open neighbourhood containing the point $y'' = (y', s'')$. Then we
obtain an affine $(Z'', z'')$ by (3). Moreover $Z_{S''} \to X_{S''}$
is a closed immersion and $Z'' \to Z_{S''}$ is an \'etale
morphism. Hence
Lemma \ref{lemma-lift-etale}
applies and we can find an \'etale morphism $X'' \to X_{S'}$ of affines
such that $Z'' \cong X'' \times_{X_{S'}} Z_{S'}$. Denote $i'' : Z'' \to X''$
the corresponding closed immersion. Setting $x'' = i''(z'')$ we obtain a
commutative diagram as in the lemma.
Properties (1), (2), (3), and (4) hold by construction.
Thus it suffices to show that (5) holds for a suitable choice of
$(S'', s'') \to (S', s')$ and $Y''$.

\medskip\noindent
We first list those properties which hold for any choice of
$(S'', s'') \to (S', s')$ and $Y''$ as in the first paragraph.
As we have $Z'' = X'' \times_X Z$ by construction we see that
$i''_*\mathcal{G}'' = \mathcal{F}''$ (with notation as in the
statement of the lemma), see
Cohomology of Schemes, Lemma \ref{coherent-lemma-affine-base-change}.
Set $n = \dim(\text{Supp}(\mathcal{F}_s)) = \dim_x(\text{Supp}(\mathcal{F}_s))$.
The morphism $Y'' \to S''$ is smooth of relative dimension $n$
(because $Y' \to S'$ is smooth of relative dimension $n$
as the composition $Y' \to Y_{S'} \to S'$ of an \'etale and
smooth morphism of relative dimension $n$ and because base change
preserves smooth morphisms of relative dimension $n$).
We have $\kappa(y'') = \kappa(y)$ and $\kappa(s) = \kappa(s'')$
hence $\kappa(y'')$ is a purely transcendental extension of $\kappa(s'')$.
The morphism of fibres $X''_{s''} \to X_s$ is an \'etale morphism of affine
schemes over $\kappa(s) = \kappa(s'')$ mapping the point $x''$ to the
point $x$ and pulling back $\mathcal{F}_s$ to $\mathcal{F}''_{s''}$.
Hence
$$
\dim(\text{Supp}(\mathcal{F}''_{s''})) =
\dim(\text{Supp}(\mathcal{F}_s)) = n =
\dim_x(\text{Supp}(\mathcal{F}_s)) =
\dim_{x''}(\text{Supp}(\mathcal{F}''_{s''}))
$$
because dimension is invariant under \'etale localization, see
Descent, Lemma \ref{descent-lemma-dimension-at-point-local}.
As $\pi'' : Z'' \to Y''$ is the base change of $\pi$ we see that
$\pi''$ is finite and as $\kappa(y) = \kappa(y'')$ we see that
$\pi^{-1}(\{y''\}) = \{z''\}$.

\medskip\noindent
At this point we have verified all the conditions of
Definition \ref{definition-one-step-devissage}
except we have not verified that $Y'' \to S''$ has geometrically
irreducible fibres. Of course in general this is not going to be
true, and it is at this point that we will use that
$\kappa(s) \subset \kappa(y)$ is purely transcendental. Namely,
let $T \subset Y'_{s'}$ be the irreducible component of
$Y'_{s'}$ containing $y' = (y, s')$. Note that $T$ is an open subscheme
of $Y'_{s'}$ as this is a smooth scheme over $\kappa(s')$. By
Varieties,
Lemma \ref{varieties-lemma-geometrically-connected-if-connected-and-point}
we see that $T$ is geometrically connected because $\kappa(s') = \kappa(s)$
is algebraically closed in $\kappa(y') = \kappa(y)$.
As $T$ is smooth we see that $T$ is geometrically irreducible. Hence
More on Morphisms,
Lemma \ref{more-morphisms-lemma-normal-morphism-irreducible}
applies and we can find an elementary \'etale morphism
$(S'', s'') \to (S', s')$ and an affine open $Y'' \subset Y'_{S''}$
such that all fibres of $Y'' \to S''$ are geometrically irreducible
and such that $T = Y''_{s''}$. After shrinking (first $Y''$ and then $S''$)
we may assume that both $Y''$ and $S''$ are affine.
This finishes the proof of the lemma.
\end{proof}

\end{document}
